\section{Introduction}
\label{sec:introduction}

Large language models increasingly delegate factual recall to an
external context: a set of retrieved documents, a cached memory
fragment, a tool-returned snippet, or a user-supplied note. The
correctness of a downstream answer therefore depends less on what the
model itself knows and more on which evidence the pipeline chooses to
place in its context window. In tasks such as medical question answering or legal reasoning, where
evidence quality matters as much as topical match, that choice has to
consider more than just relevance. A classical view of trust decomposes
into orthogonal dimensions: \emph{relevance} (does the source address
the query?), \emph{authority} (is the source intrinsically credible?),
and \emph{consistency} (do independent sources agree?). Two further
dimensions, \emph{provenance} (which specific source is this?) and
\emph{recency} (has the fact since moved?), require per-source metadata
that many standard retrieval benchmarks, including the one we study,
do not carry; we return to those two in the future-work discussion. We
refer to the three measurable dimensions collectively as \emph{trust
signals}.

Most retrieval-augmented systems, however, collapse trust to a single
number: a relevance score from a dense retriever or a sparse ranker.
This is convenient, but the single-signal view is known to fail in
predictable ways. Irrelevant-yet-topical distractors are scored
high~\cite{yoran2024makingrag}; the position of the correct passage
matters almost as much as its content~\cite{liu2024lostinthemiddle};
and conflicting passages are merged without explicit adjudication
~\cite{xie2024adaptive}. A natural response is to compose a richer
score from multiple signals. The intuition is simple: trustworthy
equals relevant plus authoritative plus corroborated. Yet this
intuition rests on an assumption that has not been tested empirically:
that each individual trust heuristic actually points toward ground
truth in the retrieval regimes we care about. In distractor-rich
settings, where candidates are deliberately topically close, two
textbook heuristics can in fact reverse sign.

We test this assumption on HotpotQA's distractor
setting~\cite{yang2018hotpotqa}. Each question is paired with ten
candidate paragraphs (two gold, eight distractors) and sentence-level
supporting-fact labels, so a paragraph-level scorer can be evaluated
directly against a provenance oracle. We define six per-paragraph
signals grouped into three trust dimensions: relevance (SBERT, BM25,
TF-IDF), authority (log-length, question-entity overlap), and
consistency (mean cross-paragraph SBERT cosine). We then contrast two
combinations. The first averages three signals uniformly, following
the textbook intuition above. The second fits a six-feature logistic
regression on only one hundred questions per seed, reserving two
hundred for evaluation.

The learned combination lifts Precision@2 to \textbf{0.345 $\pm$ 0.032}
versus SBERT's 0.312 $\pm$ 0.035, a 3.3 absolute-point improvement
that holds on every seed. More tellingly, the calibration gap between
gold and distractor scores more than doubles, from 0.170 to 0.349.
The uniform combination, despite encoding the same textbook intuition,
collapses to 0.080, well below every informed baseline including BM25
at 0.187. Inspecting the fitted weights explains both outcomes at
once: in a distractor-rich pool, paragraph length and cross-paragraph
consistency are learned as \emph{anti}-signals, because distractors
are both length-matched to gold and topically close enough to inflate
cross-source similarity above that of the true gold pair. The
classical trust heuristic is correct in prior but wrong in sign.

We make three contributions:
\begin{itemize}
  \item We formalize paragraph-level trust-aware context selection as
  a per-candidate scoring problem over six orthogonal signals covering
  relevance, authority, and consistency, and release the signals as a
  drop-in CPU-runnable feature bank.
  \item We show that a six-feature logistic regression trained on one
  hundred supervised questions per seed beats the strongest
  single-signal baseline by 3.3 Precision@2 points and more than
  doubles the gold-versus-distractor score gap, with consistent sign
  across all three seeds.
  \item We identify a previously-undiagnosed failure mode of naive
  trust-signal averaging: in distractor-rich retrieval, two of the
  three signals that textbook intuition would combine with positive
  weight are learned with \emph{negative} weight. The ablation
  isolates the phenomenon; the learned classifier recovers from it.
\end{itemize}

The rest of the paper is organized as follows. Section~\ref{sec:related-work}
surveys retrieval-augmented generation, multi-hop evidence selection,
and source attribution. Section~\ref{sec:method} formalizes the
scoring problem and the two TrustScore variants.
Section~\ref{sec:experiments} describes the protocol on HotpotQA
distractor. Section~\ref{sec:results} reports the main comparison,
the calibration diagnostic, and the learned weight pattern.
Section~\ref{sec:conclusion} discusses scope, limitations, and a
natural next step toward generation-time trust-aware selection.
